% 【本文件写什么】impl 实现层：null-stub
% - 定位：/dev/null 虚拟字符设备。
% - crate 路径：os/components/wateros-driver/driver-character/character-impl/impl-null-stub/src/
% - 如何实现对应 api-v0 trait：关键类型、状态机、数据结构
% - 算法或硬件交互流程，可用伪代码/流程图
% - 本 impl 启用的 Cargo [features] 与 arch feature，impl-riscv64 等
% - 与其它 impl 变体的差异、切换 feature 时的影响
% - 性能/内存假设与已知限制
% 【事实来源】
% - 上述 crate 源码与 Cargo.toml
% - os/feature-tree.txt 中指向本 impl 的 feature 链

\subsubsection{impl-null-stub}

\code{NullCharacterDevice}实现POSIX\code{/dev/null}语义：\code{read}返回0表示EOF；\code{write}返回\code{buf.len()}表示全部丢弃；\code{device\_kind}为\code{Null}。无MMIO、无缓冲区分配，可在任意上下文持锁调用。

\code{register\_null\_stub}与RTC stub一并由\code{register\_builtin\_character\_devices}注册，不依赖DTB节点。用户态脚本将stderr重定向到\code{/dev/null}或打开\code{O\_WRONLY}写丢弃时经VFS/devfs路由到该实现。

\begin{rustcode}
impl CharacterDevice for NullCharacterDevice {
    fn read(&mut self, _buf: &mut [u8]) -> DriverResult<usize> {
        Ok(0)
    }
    fn write(&mut self, buf: &[u8]) -> DriverResult<usize> {
        Ok(buf.len())
    }
    fn device_kind(&self) -> CharacterDeviceKind {
        CharacterDeviceKind::Null
    }
}
\end{rustcode}
